\section*{Sets and Indices}

There are no nontrivial sets or indices in the implemented model; it is a single-period truck-count problem with two truck types.

\section*{Parameters}

All parameters are scalar inputs from \texttt{general\_parameters.csv}:

\[
\begin{aligned}
R_A &:\text{ refrigerated capacity of a Type A truck } (\texttt{refrigerated\_capacity\_type\_a}),\\
N_A &:\text{ non-refrigerated capacity of a Type A truck } (\texttt{non\_refrigerated\_capacity\_type\_a}),\\
R_B &:\text{ refrigerated capacity of a Type B truck } (\texttt{refrigerated\_capacity\_type\_b}),\\
N_B &:\text{ non-refrigerated capacity of a Type B truck } (\texttt{non\_refrigerated\_capacity\_type\_b}),\\
D_R &:\text{ refrigerated cargo requirement } (\texttt{refrigerated\_cargo\_requirement}),\\
D_N &:\text{ non-refrigerated cargo requirement } (\texttt{non\_refrigerated\_cargo\_requirement}),\\
c_A &:\text{ rental cost of a Type A truck } (\texttt{rental\_cost\_type\_a}),\\
c_B &:\text{ rental cost of a Type B truck } (\texttt{rental\_cost\_type\_b}),\\
\bar A &:\text{ base fleet threshold for Type A trucks } (\texttt{base\_fleet\_type\_a}),\\
\bar B &:\text{ base fleet threshold for Type B trucks } (\texttt{base\_fleet\_type\_b}),\\
p_A &:\text{ penalty cost per Type A truck above base } (\texttt{penalty\_cost\_type\_a}),\\
p_B &:\text{ penalty cost per Type B truck above base } (\texttt{penalty\_cost\_type\_b}),\\
T_B &:\text{ Type B recovery threshold } (\texttt{type\_b\_recovery\_threshold}),\\
F_B &:\text{ fixed Type B recovery desk fee } (\texttt{type\_b\_recovery\_fee}),\\
T_F &:\text{ total fleet surge threshold } (\texttt{total\_fleet\_surge\_threshold}),\\
F_F &:\text{ fixed total fleet surge fee } (\texttt{total\_fleet\_surge\_fee}),\\
M &:\text{ big-}M\text{ constant used in the code } (=10000).
\end{aligned}
\]

For this instance, the data are:
\[
\begin{aligned}
&R_A=20,\; N_A=40,\; R_B=30,\; N_B=30,\; D_R=3000,\; D_N=4000,\\
&c_A=30,\; c_B=40,\; \bar A=30,\; \bar B=40,\; p_A=10,\; p_B=15,\\
&T_B=65,\; F_B=600,\; T_F=115,\; F_F=100,\; M=10000.
\end{aligned}
\]

\section*{Decision Variables}

\[
\begin{aligned}
x &\in \mathbb{Z}_{\ge 0} &&\text{number of Type A trucks rented (code: \texttt{x})},\\
y &\in \mathbb{Z}_{\ge 0} &&\text{number of Type B trucks rented (code: \texttt{y})},\\
x^{\text{extra}} &\in \mathbb{Z}_{\ge 0} &&\text{Type A trucks counted for base-fleet penalty (code: \texttt{x\_extra})},\\
y^{\text{extra}} &\in \mathbb{Z}_{\ge 0} &&\text{Type B trucks counted for base-fleet penalty (code: \texttt{y\_extra})},\\
z_B &\in \{0,1\} &&\text{indicator that the Type B recovery desk is opened (code: \texttt{y\_over})},\\
z_F &\in \{0,1\} &&\text{indicator that the total-fleet surge desk is opened (code: \texttt{fleet\_surge})}.
\end{aligned}
\]

\section*{Objective}

Minimize total rental, over-base penalties, and fixed recovery/surge fees:
\[
\min\;
c_A x + c_B y + p_A x^{\text{extra}} + p_B y^{\text{extra}} + F_B z_B + F_F z_F.
\]

\section*{Constraints}

Cargo requirements:
\[
R_A x + R_B y \ge D_R,
\]
\[
N_A x + N_B y \ge D_N.
\]

Base-fleet excess accounting:
\[
x^{\text{extra}} \ge x - \bar A,
\]
\[
y^{\text{extra}} \ge y - \bar B.
\]

Type B recovery-desk gate:
\[
y \le T_B + M z_B.
\]

Total-fleet surge gate:
\[
x + y \le T_F + M z_F.
\]

Integrality and binary restrictions:
\[
x,y,x^{\text{extra}},y^{\text{extra}} \in \mathbb{Z}_{\ge 0},\qquad
z_B,z_F \in \{0,1\}.
\]

\section*{Notes on Implementation}

The formulation above matches the implemented mixed-integer model in \texttt{current\_heuristic.py}.

The business rule revision is implemented as a binary gate, not as a smooth surcharge: if \(y>T_B\), the model must set \(z_B=1\) and incur \(F_B\); separately, if \(x+y>T_F\), the model must set \(z_F=1\) and incur \(F_F\). Because the constraints are only one-way big-\(M\) implications, the objective is what makes the indicators stay at \(0\) unless crossing the corresponding threshold is necessary or advantageous.

Also, the over-base penalty variables are modeled exactly as in code: \(x^{\text{extra}}\) and \(y^{\text{extra}}\) only satisfy lower bounds against \(x-\bar A\) and \(y-\bar B\), together with nonnegativity and positive objective coefficients. Thus, at optimality they take
\[
x^{\text{extra}}=\max\{0,x-\bar A\},\qquad
y^{\text{extra}}=\max\{0,y-\bar B\},
\]
but this equality is induced by minimization, not imposed directly as explicit max constraints.

The model is solved as a MIP by Gurobi through PuLP compatibility, not by exhaustive enumeration or dynamic programming.
